% GENERATED FILE. Edit canonical lessons, then run npm run generate:books.
% canonical-source-hash: fnv1a64:28d67b9918ee03be
% canonical-lessons: BN-C104-mapa, BN-C104-darano, BN-C104-shoya, BN-C104-bondho, BN-C104-mochha

\chapter{To measure, To stand, To lie down, Closed, To wipe}
\label{ch:bn-to-measure-to-stand-to-lie-down-closed-to-wipe}
\hlchaptermodality{104}

\begin{chapteropening}
\textbf{By the end of this chapter:} \emph{I can say to measure, to stand, to lie down, closed and to wipe.}

Complete the last lesson of chapter 104: i can say to measure, to stand, to lie down, closed and to wipe.
\end{chapteropening}

\section[\texorpdfstring{\bn{মাপা} (māpā)}{মাপা (māpā)}]{\bn{মাপা} (māpā) — to measure}

\label{lesson:BN-C104-mapa}



\begin{quote}
Before the new one: say the Bengali for to hit, then the Bengali for a share.
\end{quote}

\subsection*{You'll want to know: \bn{মাপা}}
\textbf{\bn{মাপা}} — \emph{māpā} — \textquotedblleft{}to measure\textquotedblright{}.

\begin{grammarlens}[title={What you've built}]
One more word for this chapter.
\end{grammarlens}

\subsection*{Your turn}
\begin{itemize}
\raggedright
  \item \emph{Say it:} \emph{māpā}
  \item \emph{Say it:} \emph{māpā}, once more
  \item \emph{Say it:} say \emph{māpā}
  \item \emph{Recall:} say the Bengali for to hit, then the Bengali for a share, then say \emph{māpā} again
\end{itemize}

\begin{tcolorbox}[breakable,colback=teal!4,colframe=teal!35!black,title={Before you move on}]
What does \bn{মাপা} mean? (\textquotedblleft{}To measure\textquotedblright{}.) Say it once more.
\end{tcolorbox}
\section[\texorpdfstring{\bn{দাঁড়ানো} (dā̃ṛāno)}{দাঁড়ানো (dā̃ṛāno)}]{\bn{দাঁড়ানো} (dā̃ṛāno) — to stand}

\label{lesson:BN-C104-darano}



\begin{quote}
Before the new one: say the Bengali for a share, then the Bengali for to measure.
\end{quote}

\subsection*{You'll want to know: \bn{দাঁড়ানো}}
\textbf{\bn{দাঁড়ানো}} — \emph{dā̃ṛāno} — \textquotedblleft{}to stand\textquotedblright{}.

\begin{grammarlens}[title={What you've built}]
One more word for this chapter.
\end{grammarlens}

\subsection*{Your turn}
\begin{itemize}
\raggedright
  \item \emph{Say it:} \emph{dā̃ṛāno}
  \item \emph{Say it:} \emph{dā̃ṛāno}, once more
  \item \emph{Say it:} say \emph{dā̃ṛāno}
  \item \emph{Recall:} say the Bengali for a share, then the Bengali for to measure, then say \emph{dā̃ṛāno} again
\end{itemize}

\begin{tcolorbox}[breakable,colback=teal!4,colframe=teal!35!black,title={Before you move on}]
What does \bn{দাঁড়ানো} mean? (\textquotedblleft{}To stand\textquotedblright{}.) Say it once more.
\end{tcolorbox}
\section[\texorpdfstring{\bn{শোয়া} (shoyā)}{শোয়া (shoyā)}]{\bn{শোয়া} (shoyā) — to lie down}

\label{lesson:BN-C104-shoya}



\begin{quote}
Before the new one: say the Bengali for to measure, then the Bengali for to stand.
\end{quote}

\subsection*{You'll want to know: \bn{শোয়া}}
\textbf{\bn{শোয়া}} — \emph{shoyā} — \textquotedblleft{}to lie down\textquotedblright{}.

\begin{grammarlens}[title={What you've built}]
One more word for this chapter.
\end{grammarlens}

\subsection*{Your turn}
\begin{itemize}
\raggedright
  \item \emph{Say it:} \emph{shoyā}
  \item \emph{Say it:} \emph{shoyā}, once more
  \item \emph{Say it:} say \emph{shoyā}
  \item \emph{Recall:} say the Bengali for to measure, then the Bengali for to stand, then say \emph{shoyā} again
\end{itemize}

\begin{tcolorbox}[breakable,colback=teal!4,colframe=teal!35!black,title={Before you move on}]
What does \bn{শোয়া} mean? (\textquotedblleft{}To lie down\textquotedblright{}.) Say it once more.
\end{tcolorbox}
\section[\texorpdfstring{\bn{বন্ধ} (bôndhô)}{বন্ধ (bôndhô)}]{\bn{বন্ধ} (bôndhô) — closed, shut}

\label{lesson:BN-C104-bondho}



\begin{quote}
Before the new one: say the Bengali for to stand, then the Bengali for to lie down.
\end{quote}

\subsection*{You'll want to know: \bn{বন্ধ}}
\textbf{\bn{বন্ধ}} — \emph{bôndhô} — \textquotedblleft{}closed, shut\textquotedblright{}.

\begin{grammarlens}[title={What you've built}]
One more word for this chapter.
\end{grammarlens}

\subsection*{Your turn}
\begin{itemize}
\raggedright
  \item \emph{Say it:} \emph{bôndhô}
  \item \emph{Say it:} \emph{bôndhô}, once more
  \item \emph{Say it:} say \emph{bôndhô}
  \item \emph{Recall:} say the Bengali for to stand, then the Bengali for to lie down, then say \emph{bôndhô} again
\end{itemize}

\begin{tcolorbox}[breakable,colback=teal!4,colframe=teal!35!black,title={Before you move on}]
What does \bn{বন্ধ} mean? (\textquotedblleft{}Closed, shut\textquotedblright{}.) Say it once more.
\end{tcolorbox}
\section[\texorpdfstring{\bn{মোছা} (mochhā)}{মোছা (mochhā)}]{\bn{মোছা} (mochhā) — to wipe}

\label{lesson:BN-C104-mochha}



\begin{quote}
Before the new one: say the Bengali for to lie down, then the Bengali for closed.
\end{quote}

\subsection*{You'll want to know: \bn{মোছা}}
\textbf{\bn{মোছা}} — \emph{mochhā} — \textquotedblleft{}to wipe\textquotedblright{}.

\begin{grammarlens}[title={What you've built}]
One more word for this chapter.
\end{grammarlens}

\subsection*{Your turn}
\begin{itemize}
\raggedright
  \item \emph{Say it:} \emph{mochhā}
  \item \emph{Say it:} \emph{mochhā}, once more
  \item \emph{Say it:} say \emph{mochhā}
  \item \emph{Recall:} say the Bengali for to lie down, then the Bengali for closed, then say \emph{mochhā} again
\end{itemize}

\begin{tcolorbox}[breakable,colback=teal!4,colframe=teal!35!black,title={Before you move on}]
What does \bn{মোছা} mean? (\textquotedblleft{}To wipe\textquotedblright{}.) Say it once more.
\end{tcolorbox}
